\section{The four-seam carrier and the transfer}
\label{sec:carrier}

\subsection{Incidence graph and the schedule}

It is convenient to split each physical point $x$ into two \emph{incidences}
$x^{P}$ and $x^{Q}$ joined by a \emph{vertex edge}; an arch of $P$ joins the
$P$-incidences of its endpoints, and likewise for $Q$. The resulting graph is a
subdivision-like replacement of the overlay: it is $2$-regular and connected
exactly when the overlay is.

Fix a \HIGH\ sector $(L,u,v)$, and write $R=2n-L$, $\Delta=|L-n|$,
$d=n-K$. The \emph{long} side is the longer of the two halves, of length
$n+\Delta$, and the \emph{short} side has length $n-\Delta$ (at $\Delta=0$ take
the left side as long). The \emph{schedule} processes physical points inward:
first $\Delta$ points of the long side, then $n-\Delta$ pairs consisting of one
long-side and one short-side point, then the remaining $\Delta$ long-side
points. Each side is thus read in its own inward order, and every point is
visited exactly once. The \LOW\ branch uses the chronological schedule
($L=2n$, $R=0$).

Both matchings are cut at $L$; their halves have budgets
\eqref{eq:budgets} with $w=u$ for $P$ and $w=v$ for $Q$.
A \emph{source word} of the sector consists of two inward ballot words for $P$
and two for $Q$ with these lengths and budgets \emph{and} whose two left words
satisfy the first-hit condition of \cref{def:sectors}: the original grade
$j-c_P(j)-c_Q(j)$ is $<K$ after $j<L$ left points and $=K$ after $L$. Lengths
and budgets alone do not imply this; a pair may have the prescribed heights at
$L$ but reach $K$ earlier, and then it belongs to another sector. A source
word of \LOW\ consists of two complete ballot words of length $2n$ ending at
height $0$ whose grade is $<K$ after every prefix. We call these the
\emph{sector-valid} source words; by \cref{lem:cut} they are in bijection with
the pairs $(P,Q)$ of the sector. Reading one point of the
schedule reveals one letter of $P$ and one of $Q$ on that side; the letter pair
is one of the four \emph{labels} $UU,UD,DU,DD$.

\subsection{The state}

At any stage of the schedule the processed graph consists of the vertex edges
of processed points, the internal arches determined by the stack matches so
far, and the through arches already joined by \cref{cor:fifo}. Its nodes have
degree $2$ (saturated) or $1$ (a \emph{port}). Live components are paths whose
two ports are their endpoints, plus possibly closed cycles.

The ports are exactly the unmatched, not yet joined openings of the four
half words. Group them by owner and side into $PL$, $PR$, $QL$, $QR$, each
ordered oldest first. For a group $g$ on side $s$ with $i_s$ processed points,
$c_g$ letters $D$ read, budget $A_g$ and opposite group $g'$ (same owner, other
side), \cref{lem:forced} gives
\begin{equation}
\label{eq:counters}
 h_g=i_s-2c_g,\quad f_g=\max(0,\,i_s-A_g-c_g),\quad
 \rho_g=\min(h_g,A_g-c_g),\quad q_g=\max(f_g-f_{g'},0),
\end{equation}
where $h_g$ is the original height, $f_g$ the cumulative emission (written $f$ rather than $E$, which is reserved for the observer channel of \cref{sec:open}), and the
retained group is a \emph{FIFO} of $q_g$ emitted-but-unjoined openings (oldest
first) followed by a \emph{reservoir} of $\rho_g$ openings. The number of
retained ports of one owner is
\begin{equation}
\label{eq:ports}
 p=h_{L}+h_{R}-2\min(f_L,f_R).
\end{equation}

\begingroup
% Slightly more interword shrink keeps the counter tuple on one line.
\spaceskip=3.33333pt plus 1.66667pt minus 1.33333pt
\begin{definition}[Carrier key]
\label{def:key}
\lean{Meanders.FirstCrossing.Key}
\leanok
\uses{cor:fifo, lem:forced}
The \emph{key} of a stage is the tuple of four down counters
$(c_{PL},c_{PR},c_{QL},c_{QR})$ together with the pairing of the retained
ports into path endpoints, recorded as a fixed-point-free involution on the
positions of the cyclic sequence
\[
 \rev(PL),\ PR,\ \rev(QR),\ QL .
\]
The sector metadata and the stage index are shared by a whole layer and are
not part of the key.
\end{definition}
\endgroup

Given the stage, the four counters determine every group length and every
FIFO/reservoir separator by \eqref{eq:counters}. No component names, ages,
path lengths or history are stored.

\begin{figure}[t]
\centering
\begin{tikzpicture}[x=0.9cm,y=0.9cm,>=Stealth,every node/.style={font=\small}]
 % line
 \draw[thick] (0,0) -- (12,0);
 % processed blocks
 \fill[gray!15] (0,-2.6) rectangle (3.6,2.6);
 \fill[gray!15] (8.4,-2.6) rectangle (12,2.6);
 \node at (1.8,2.9) {processed left prefix ($i_L\le L$ points)};
 \node at (10.2,2.9) {processed right prefix ($i_R\le R$ points)};
 \node at (6,-2.45) {unprocessed};
 \draw[dashed] (3.6,-2.6) -- (3.6,2.6);
 \draw[dashed] (8.4,-2.6) -- (8.4,2.6);
 % PL arcs (upper, left): pending, crossing the left cut
 \draw (0.6,0) arc (180:90:2.4 and 2.2) -- (3.9,2.2);
 \draw (1.4,0) arc (180:90:1.6 and 1.5) -- (3.9,1.5);
 \draw (2.2,0) arc (180:90:0.8 and 0.8) -- (3.9,0.8);
 % internal PL arch
 \draw (2.6,0) arc (180:0:0.35 and 0.35);
 % QL arcs (lower, left)
 \draw (1.0,0) arc (180:270:2.0 and 1.8) -- (3.9,-1.8);
 \draw (1.8,0) arc (180:270:1.2 and 1.0) -- (3.9,-1.0);
 % PR arcs (upper, right)
 \draw (11.4,0) arc (0:90:2.4 and 2.0) -- (8.1,2.0);
 \draw (10.6,0) arc (0:90:1.6 and 1.2) -- (8.1,1.2);
 % QR arcs (lower, right)
 \draw (11.0,0) arc (0:-90:2.0 and 1.9) -- (8.1,-1.9);
 \draw (10.2,0) arc (0:-90:1.2 and 1.2) -- (8.1,-1.2);
 \draw (9.4,0) arc (0:-90:0.5 and 0.5) -- (8.1,-0.5);
 % group labels
 \node[anchor=east,fill=white,inner sep=1pt] at (3.5,1.5) {$PL$};
 \node[anchor=east,fill=white,inner sep=1pt] at (3.5,-1.4) {$QL$};
 \node[anchor=west,fill=white,inner sep=1pt] at (8.5,1.6) {$PR$};
 \node[anchor=west,fill=white,inner sep=1pt] at (8.5,-1.2) {$QR$};
 % seams
 \node[draw,fill=white,rounded corners=1pt,inner sep=1.5pt] at (3.6,0) {left vertex seam};
 \node[draw,fill=white,rounded corners=1pt,inner sep=1.5pt] at (8.4,0) {right vertex seam};
 \node[draw,fill=white,rounded corners=1pt,inner sep=1.5pt] at (6,2.2) {$P$ through seam};
 \node[draw,fill=white,rounded corners=1pt,inner sep=1.5pt] at (6,-2.0) {$Q$ through seam};
 % orientation of cyclic order
 \draw[->,gray] (4.3,1.9) -- (4.3,0.4);
 \node[gray,anchor=west] at (4.35,1.15) {\scriptsize newest};
\end{tikzpicture}
\caption{The four-seam carrier at an intermediate stage; the shaded blocks
are the processed prefixes of the two halves, not the complete halves, and
the arches are schematic. Pending arches cross the two cuts; walking
around the unprocessed middle region meets the ports in the cyclic order
$\rev(PL),PR,\rev(QR),QL$. The newest openings of each group are nearest the
corresponding cut; the oldest are farthest from it. The two vertex seams sit on the line at
the cuts; the two through seams sit between the oldest ports of the two halves
of each owner.}
\label{fig:carrier}
\end{figure}

\subsection{Transitions}

Processing a point on side $s$ with label $(\alpha,\beta)$, $\alpha$ for $P$
and $\beta$ for $Q$, consists of the following steps. Every operation is
performed on the current key and produces a successor key or a rejection.
\begin{enumerate}[nosep,label=(\arabic*)]
\item \emph{Feasibility.} A $D$ letter needs a nonempty reservoir of its group
  and remaining budget $A_g-c_g\ge1$; a $U$ letter needs enough remaining
  length to spend the remaining budget. The grade after this point,
  computed from the original counters (\cref{lem:filter}), must satisfy the
  sector's first-hit condition: in \LOW\ it must be $<K$ after every point;
  on the left side of \HIGH\ it must be $<K$ after the $j$-th left point for
  $j<L$ and $=K$ for $j=L$. Right-side points impose no grade test.
\item \emph{Vertex surgery at the vertex seam of side $s$.} The new point
  contributes two adjacent ports $x^{P},x^{Q}$ joined by a vertex edge. Label
  $UU$ inserts this \emph{cup} at the seam. Label $UD$ (resp.\ $DU$) attaches
  the $Q$ (resp.\ $P$) incidence to the newest port of its reservoir, a
  \emph{replacement} of that port by $x^{P}$ (resp.\ $x^{Q}$) at the seam.
  Label $DD$ attaches both incidences to the two newest reservoir ports of $P$
  and $Q$, which are adjacent across the seam: a \emph{cap}. A cap either
  splices two paths into one or, if its two ports were endpoints of the same
  path, closes a cycle.
\item \emph{Emission.} Update the counters and recompute $f_g$ by
  \eqref{eq:counters}; a newly emitted opening moves from the reservoir to the
  FIFO by shifting a separator.
\item \emph{Forced joins.} While both FIFOs of an owner are nonempty, cap their
  two oldest ports, which are adjacent at that owner's through seam; first for
  $P$, then for $Q$. Each point causes at most one join per owner.
\item \emph{Cycle rule.} Reject if a cycle closed in steps (2) or (4), unless
  this is the last point of the schedule, exactly one cycle closed, and no port
  remains; in that case the state is \emph{accepted}.
\item \emph{Canonicalization.} Rename the retained ports by their positions in
  the four groups.
\end{enumerate}

\begin{lemma}[Four-seam invariant]
\label{lem:seams}
\lean{Meanders.FirstCrossing.adjacentCup_representsPaths, Meanders.FirstCrossing.adjacentCap_representsPaths, Meanders.FirstCrossing.cyclicCheck_iff_noncrossing}
\leanok
\uses{def:key, cor:fifo}
At every stage the pairing of the retained ports into path endpoints is
noncrossing with respect to the cyclic order $\rev(PL),PR,\rev(QR),QL$, and every
operation in steps (2) and (4) acts on two adjacent positions of that order.
\end{lemma}

\begin{proof}
\leanok
A point on the left side is adjacent, in the drawing, to the newest opening of
$PL$ above the line and the newest opening of $QL$ below it; a point on the
right side is adjacent to the newest openings of $PR$ and $QR$. The oldest
openings of $PL$ and $PR$ are the outermost pending arches above the line and
are joined across the middle region; likewise for $QR$ and $QL$ below it. Thus
the four seams are the four places in the cyclic sequence where a cup,
replacement or cap can occur (\cref{fig:carrier}). Starting from the empty
boundary, an adjacent cup adds a pair of neighbouring endpoints, a replacement
moves one endpoint across an empty seam without changing its order relative to
the others, and an adjacent cap either splices two paths (the surviving two
endpoints remain paired with each other in the induced order) or closes a
cycle. Each preserves noncrossingness. No relation between the two processed
lengths is used, so the argument covers every stage of the asymmetric schedule.
\end{proof}

\begin{lemma}[Chronological filter]
\label{lem:filter}
\lean{Meanders.FirstCrossing.firstHeightCounters, Meanders.FirstCrossing.Source.prefix_allowed}
\leanok
\uses{def:sectors, def:key}
Let $i_L$ points of the left side have been processed. Then the original grade
after $i_L$ points equals $i_L-c_{PL}-c_{QL}$, whatever has been processed on
the right side and whatever joins have been performed. Testing this quantity
after each left point implements \cref{def:sectors} exactly.
\end{lemma}

\begin{proof}
\leanok
The counters record letters of the original words and are not changed by joins
or by right-side processing, so $h_P(i_L)=i_L-2c_{PL}$ and likewise for $Q$.
The first-hit condition is a condition on the prefix heights only.
\end{proof}

\begin{lemma}[Cycle rule]
\label{lem:cycle}
\lean{Meanders.FirstCrossing.saturatedComponent_permanent, Meanders.FirstCrossing.acceptedSourceEquiv, Meanders.FirstCrossing.CoversProcessed.connected_of_one_cycle}
\leanok
\uses{lem:cut, cor:fifo, lem:seams}
A sector-valid source word (\cref{sec:carrier}, first subsection) is accepted
by the transfer if and only if its matching pair has connected overlay. A
word with the prescribed lengths and budgets that violates the first-hit
condition is rejected by step~(1) whatever its connectivity.
\end{lemma}

\begin{proof}
\leanok
For a sector-valid word the grade tests of step~(1) pass at every left point
(\cref{lem:filter}), so acceptance is decided by the cycle rule alone. Every
edge inserted by the transfer is an edge of the final incidence graph
(\cref{lem:cut,cor:fifo}). If a cycle closes before the last point, its nodes
all have degree $2$ and no later edge can attach to it; since the last point
still adds a vertex edge outside it, the final graph is disconnected. If the
final graph is connected it is a single cycle, and any cycle of a subgraph of a
cycle is the whole cycle, so the only closure happens at the last operation, and
then no port remains. Conversely a final state with exactly one closure and no
ports is connected. Since the schedule is deterministic and the cut bijection
recovers the source from its four half words, each source word has exactly
one label history, so accepted histories and connected pairs correspond
bijectively.
\end{proof}

\begin{lemma}[Future equivalence]
\label{lem:future}
\lean{Meanders.FirstCrossing.stateFutureEquiv}
\leanok
\uses{def:key, lem:filter, lem:cycle, lem:seams}
Two partial source words of the same sector that reach the same stage with the
same key have the same set of legal labelled continuations, with the same
acceptance decisions.
\end{lemma}

\begin{proof}
\leanok
The stage and the counters determine all group lengths, separators, budgets and
reservoir tops by \eqref{eq:counters}, so every feasibility test in step (1),
the positions used in steps (2) and (4), the emissions in step (3) and the
filter of \cref{lem:filter} coincide. The pairing determines whether a cap
splices or closes, hence the cycle rule. The successors again share a key, and
induction on the remaining schedule finishes the proof.
\end{proof}

\subsection{The transfer}

For each sector the algorithm maintains a \emph{layer}: a list of distinct
keys sorted lexicographically, each with an integer weight $C$ (and a second
integer $E$ for the open observer of \cref{sec:open}). The initial layer is the
empty key with weight $1$. For each point of the schedule, every key emits up
to four successors, one per label; the successors are sorted and equal keys
are merged by adding weights. After the last point, the accepted keys contribute their weight to the sector
total; there is at most one such key, the \emph{terminal key}, whose four
counters equal the four down budgets and whose pairing is empty (it is not the
initial empty key, whose counters are zero). By
\cref{lem:future} the weight of a key equals the number of partial source
words reaching it, so the sector total is the number of connected pairs in the
sector; by \cref{lem:partition} the sum over \LOW\ and all admissible sectors is
$\Closed(n)$. Sectors are processed one after another; two logical layers are alive at
any time, together with the $O(M)$ working storage of the $M$ emitted records
and their sort, where $M$ is at most four times the current key count.
